\chapter{Revisão de Literatura}

\section{Pix, pagamentos instantâneos e risco de fraude}

O Pix é uma infraestrutura de pagamentos instantâneos operada no âmbito do Sistema de Pagamentos Brasileiro. Sua disponibilidade contínua e ampla adoção tornam relevantes mecanismos capazes de reconhecer sinais de risco com baixa latência. Em 2024, o instrumento apresentou crescimento de 52\% em quantidade de transações e respondeu por 47\% dos pagamentos sem uso de espécie no último trimestre do ano \cite{bcb_pagamentos_2024}. Esses números caracterizam escala e adoção, mas não devem ser interpretados como medida da incidência de fraude.

O risco associado a uma transferência não decorre apenas do funcionamento técnico do instrumento. Manipulação do usuário, comprometimento de credenciais, abuso de conta e mudança de comportamento podem produzir operações aparentemente regulares na camada transacional. Consequentemente, um classificador não observa diretamente a intenção do agente nem substitui a apuração de uma instituição financeira. Sua saída deve ser apresentada como estimativa de risco sujeita a revisão.

O Mecanismo Especial de Devolução foi introduzido no Regulamento do Pix pela Resolução BCB nº 103/2021 \cite{bcb_resolucao103_2021}. O mecanismo compõe um processo regulatório e operacional; não equivale a um rótulo de aprendizado de máquina, a uma garantia automática de devolução ou a uma conclusão jurídica de fraude. Como o Regulamento do Pix e o Guia do MED são atualizados, afirmações sobre procedimentos e prazos exigem versão e data de vigência \cite{bcb_resolucao1_2020,bcb_guia_med_v43}.

A Pesquisa Febraban de Tecnologia Bancária 2024 registra a presença de detecção de fraude e lavagem de dinheiro entre casos de uso de inteligência artificial declarados pelos bancos participantes \cite{febraban2024}. Trata-se de evidência setorial de pertinência, não de uma taxa de fraude Pix nem de validação de um algoritmo específico. Essa distinção é necessária para evitar que estatísticas de investimento ou adoção sejam transformadas em evidência de desempenho.

\section{Aprendizado de máquina para detecção de fraude}

Detecção de fraude financeira é um problema de classificação rara, dinâmica e sensível ao protocolo de avaliação. Entre as dificuldades recorrentes estão desbalanceamento, grande volume, adaptação dos fraudadores, atraso na confirmação do rótulo e custos assimétricos de erro \cite{abdallah2016survey,dalpozzolo2018realistic}. A ordem e o histórico das transações também carregam informação, o que motiva atributos de recência, frequência e desvio calculados somente com eventos anteriores \cite{jurgovsky2018sequence}.

Em conjuntos fortemente desbalanceados, a acurácia pode permanecer alta mesmo quando o classificador falha na classe positiva. Curvas e áreas de precisão--revocação são particularmente informativas nesse cenário \cite{saito2015precision}. Por isso, este trabalho prevê AUC-PR, AUC-ROC, precisão, revocação, F1 e matriz de confusão, acompanhadas do limiar usado. O limiar não deve ser escolhido no conjunto de teste.

A regressão logística será usada como linha de base por oferecer um ponto de comparação simples e reproduzível. Random Forest representa uma alternativa não linear baseada na agregação de árvores \cite{breiman2001random}, enquanto XGBoost usa boosting de árvores com mecanismos de regularização e escalabilidade \cite{chen2016xgboost}. Comparar esses modelos permite verificar se ganhos de complexidade se traduzem em melhoria relevante para a classe rara.

SMOTE cria amostras sintéticas da classe minoritária a partir de vizinhos no espaço de atributos \cite{chawla2002smote}. A técnica não deve ser aplicada antes da divisão dos dados: nesse caso, exemplos derivados da validação ou do teste podem contaminar o treino. A comparação adequada coloca reamostragem dentro do pipeline ajustado somente em treino e a contrasta com alternativas como ponderação de classe. Variáveis categóricas, valores ausentes e atributos de alta cardinalidade exigem tratamento compatível antes de qualquer interpolação.

Mudança de conceito limita a interpretação de uma divisão aleatória como evidência operacional. Um modelo pode reconhecer padrões presentes simultaneamente em treino e teste sem manter desempenho em período posterior. Assim, uma separação temporal é metodologicamente preferível quando a ordem relativa está disponível, ainda que \texttt{TransactionDT} não revele a data civil real. O IEEE-CIS será usado como proxy técnico de comércio eletrônico/cartão \cite{ieee_cis_2019}; seus resultados não demonstram generalização para Pix.

\section{Explicabilidade por SHAP}

SHAP expressa uma previsão como valor de referência acrescido das contribuições atribuídas aos atributos \cite{lundberg2017unified}. A análise local ajuda a identificar quais entradas elevaram ou reduziram o escore em uma transação, enquanto agregações de valores absolutos apoiam uma visão global do comportamento do modelo.

Essas contribuições não estabelecem causalidade, intenção criminosa ou significado semântico de uma variável anônima. No IEEE-CIS, grupos como \texttt{C*}, \texttt{D*}, \texttt{M*}, \texttt{V*} e \texttt{id\_*} possuem semântica individual parcial ou não publicada. Uma explicação correta pode informar que \texttt{V258} influenciou o modelo, mas não pode renomeá-la como conta, dispositivo ou destinatário Pix sem uma fonte que sustente essa equivalência.

\section{Retrieval-Augmented Generation e modelos de linguagem}

Modelos de linguagem organizam texto a partir de padrões aprendidos, mas não garantem atualização, rastreabilidade ou fidelidade factual. Retrieval-Augmented Generation combina recuperação de documentos com geração condicionada ao contexto recuperado \cite{lewis2020rag}. Em vez de depender apenas da memória paramétrica do modelo, o sistema consulta uma coleção externa e inclui trechos relevantes na entrada da geração.

Um pipeline de recuperação textual contém documentos, segmentação em trechos, embeddings, índice e retriever. Sentence-BERT mostrou uma forma eficiente de produzir representações de sentenças adequadas à comparação semântica \cite{reimers2019sentencebert}, enquanto recuperação densa representa consultas e passagens em espaços vetoriais compatíveis \cite{karpukhin2020dense}. FAISS oferece algoritmos e estruturas para busca eficiente por similaridade vetorial \cite{johnson2021billion}.

O retriever recebe uma consulta e devolve trechos candidatos; o modelo de linguagem recebe esses trechos e produz a resposta. O banco vetorial, contudo, não substitui o documento original como fonte de verdade. Cada trecho precisa preservar identificador, emissor, versão, vigência, página ou seção e URL. No domínio regulatório, misturar versões históricas e vigentes pode produzir uma resposta linguisticamente plausível e operacionalmente incorreta.

RAG também não elimina alucinação. A recuperação pode retornar material irrelevante, e o gerador pode produzir afirmações não sustentadas pelo contexto. A avaliação deve separar qualidade da recuperação, fidelidade ao documento, coerência com a evidência SHAP e clareza humana. Métricas automatizadas, como as propostas pelo RAGAS, podem apoiar essa análise, mas não substituem um conjunto de referência e revisão humana \cite{es2024ragas}.

\section{Síntese e lacuna investigada}

A literatura fornece métodos para classificação rara, explicação local e recuperação documental, porém sua simples justaposição não garante uma explicação fiel. O classificador responde quanto o padrão se aproxima da classe aprendida; SHAP registra quais entradas influenciaram a previsão; o RAG recupera contexto verificável. A lacuna investigada é integrar essas evidências sem transformar correlação em causalidade, sem inventar semântica para atributos anonimizados e sem confundir risco estatístico com decisão regulatória.

O protótipo adotará um registro versionado de atributos explicáveis. Apenas features explícitas ou derivadas com fórmula, janela, unidade e limitação aprovadas poderão fornecer conceitos ao retriever. Colunas anônimas podem permanecer úteis ao classificador, mas serão descritas como atributos anonimizados quando aparecerem na explicação.
